\documentclass[10pt,a4paper]{article}
\usepackage[a4paper,left=18mm,right=18mm,top=21mm,bottom=22mm,headheight=14pt,headsep=7mm,footskip=10mm]{geometry}
\usepackage{amsmath,amsthm}
\usepackage{txfonts}            % Times-like math, close to the paper's pslatex + newtxmath
\usepackage{latexsym}
\usepackage[mathscr]{euscript}  % as in paper.tex
\usepackage{mathtools}
\usepackage{array}
\usepackage[no-math]{fontspec}
\setmainfont{FreeSerif}[
  Path=@FREEFONT_DIR@, Extension=.ttf,
  UprightFont=*, ItalicFont=*Italic, BoldFont=*Bold, BoldItalicFont=*BoldItalic,
  Ligatures=TeXOff]
\setsansfont{Inter}[
  Path=@INTER_DIR@, Extension=.otf,
  UprightFont=*-Regular, BoldFont=*-SemiBold, ItalicFont=*-Italic, BoldItalicFont=*-SemiBoldItalic,
  Scale=0.88, Ligatures=TeXOff, RawFeature={-calt;-case}]   % contextual alternates map to private-use code points in the text layer
\setmonofont{DejaVuSansMono}[
  Path=@DEJAVU_DIR@, Extension=.ttf,
  UprightFont=*, BoldFont=*-Bold, ItalicFont=*-Oblique, BoldItalicFont=*-BoldOblique,
  Scale=0.84]
% font directories (@…@) are filled in by build.py from FREEFONT_DIR, DEJAVU_DIR, INTER_DIR
% per-character fallback fonts (XeTeX has no automatic font fallback)
\newfontfamily\FBfs{FreeSerif}[Path=@FREEFONT_DIR@, Extension=.ttf,
  UprightFont=*, ItalicFont=*Italic, BoldFont=*Bold, BoldItalicFont=*BoldItalic]
\newfontfamily\FBfm{FreeMono}[Path=@FREEFONT_DIR@, Extension=.ttf,
  UprightFont=*, ItalicFont=*Oblique, BoldFont=*Bold, BoldItalicFont=*BoldOblique, Scale=0.84]
\newfontfamily\FBdvs{DejaVuSans}[Path=@DEJAVU_DIR@, Extension=.ttf,
  UprightFont=*, ItalicFont=*-Oblique, BoldFont=*-Bold, BoldItalicFont=*-BoldOblique, Scale=0.82]
\usepackage{newunicodechar}
\input{fallback.tex}

\usepackage{xcolor}
\definecolor{ink}{HTML}{1F2328}
\definecolor{muted}{HTML}{6B7280}
\definecolor{faint}{HTML}{9CA3AF}
\definecolor{accent}{HTML}{1E3A5F}
\definecolor{linkc}{HTML}{2155A6}
\definecolor{paperbg}{HTML}{FBF7EE}
\definecolor{paperrule}{HTML}{B8892E}
\definecolor{paperlab}{HTML}{7C5512}
\definecolor{leanbg}{HTML}{F1F5FB}
\definecolor{leanrule}{HTML}{2F5D9A}
\definecolor{defbg}{HTML}{F7F7F4}
\definecolor{defframe}{HTML}{D6D6CF}
\definecolor{miscbg}{HTML}{F4F4F2}
\definecolor{synkw}{HTML}{1E3A8A}
\definecolor{synlog}{HTML}{9A3412}
\definecolor{syntype}{HTML}{0F6466}
\definecolor{synproof}{HTML}{7A808A}
\definecolor{syncomment}{HTML}{6B7280}
\definecolor{chipgreenbg}{HTML}{E6F2EA}\definecolor{chipgreen}{HTML}{1D6B3A}
\definecolor{chipbluebg}{HTML}{E8EEF7}\definecolor{chipblue}{HTML}{294C7A}
\definecolor{chipamberbg}{HTML}{FCEFD3}\definecolor{chipamber}{HTML}{8A5300}
\color{ink}

\usepackage{fvextra}
\usepackage[most]{tcolorbox}
\usepackage{enumitem}
\usepackage{needspace}
\usepackage{fancyhdr}
\usepackage{lastpage}
\usepackage{etoolbox}
\usepackage{microtype}
\usepackage[newcommands]{ragged2e}
\setlength{\RaggedRightRightskip}{0pt plus 1fil}
\setlength{\RaggedRightParindent}{0pt}
\setlength{\JustifyingParindent}{0pt}
\AtBeginDocument{\RaggedRight}
\hyphenpenalty=10000 \exhyphenpenalty=10000   % never hyphenate identifiers in prose
\usepackage[hidelinks,unicode,bookmarksnumbered=false,pdfusetitle]{hyperref}
\usepackage{bookmark}
\hypersetup{colorlinks=true,linkcolor=linkc,urlcolor=linkc,citecolor=linkc,
  pdfauthor={Rendered from PaperFormalization/Statements.lean},
  pdftitle={Statements.lean: the paper statements and their Lean encodings},
  bookmarksopen=true,bookmarksopenlevel=1}

\setlength{\parindent}{0pt}
\setlength{\parskip}{5pt plus 1.5pt minus 1pt}
\linespread{1.06}
\emergencystretch=3em
\tolerance=2000
\hbadness=3000
\raggedbottom
\setlist{itemsep=2pt,topsep=3pt,parsep=1pt,leftmargin=1.4em}
\setlist[itemize]{label={\color{muted}\textbullet}}

% ---------- break opportunity inside code spans ----------
\newcommand\BO{\penalty300\relax}
\newcommand\nohy[1]{{\language=255 #1}}
\DeclareRobustCommand\lc[1]{\texttt{#1}}
\DeclareRobustCommand\lclink[2]{\hyperlink{#1}{\texttt{#2}}}

% ---------- Lean syntax classes ----------
\newcommand\LK[1]{\textcolor{synkw}{\textbf{#1}}}
\newcommand\LL[1]{\textcolor{synlog}{#1}}
\newcommand\LT[1]{\textcolor{syntype}{#1}}
\newcommand\LN[1]{#1}
\newcommand\LC[1]{\textcolor{syncomment}{\textit{#1}}}
\newcommand\LP[1]{\textcolor{synproof}{#1}}
\newcommand\LD[2]{\hyperlink{#1}{\textcolor{linkc}{#2}}}
\newcommand\LBS{\char"5C\relax}
\newcommand\LLB{\char"7B\relax}
\newcommand\LRB{\char"7D\relax}

\fvset{fontsize=\normalsize}
\DefineVerbatimEnvironment{LeanVerb}{Verbatim}{%
  commandchars=\\\{\},breaklines=true,breakautoindent=true,breakindent=1.5em,
  breaksymbolleft={\textcolor{paperrule}{\ttfamily\scriptsize\symbol{"21AA}}},
  breaksymbolsepleft=2pt,
  numbers=left,numbersep=16pt,
  formatcom={\color{ink}},
  numberblanklines=true}
\renewcommand\theFancyVerbLine{\textcolor{faint}{\sffamily\fontsize{6.2}{7}\selectfont\arabic{FancyVerbLine}}}
\DefineVerbatimEnvironment{TexVerb}{Verbatim}{%
  breaklines=true,breakanywhere=true,breakautoindent=false,
  breaksymbolleft={\textcolor{paperrule}{\ttfamily\scriptsize\symbol{"21AA}}},
  numbers=left,numbersep=16pt,fontsize=\small}
\DefineVerbatimEnvironment{ProseVerb}{Verbatim}{%
  xleftmargin=1.8em,fontsize=\normalsize,breaklines=true}

% ---------- boxes ----------
\newcommand\boxlabel[3]{% color, label, meta
  {\sffamily\fontsize{7.2}{9}\selectfont\color{#1}\textbf{\MakeUppercase{#2}}%
   \ifstrempty{#3}{}{\hspace{0.6em}{\mdseries\color{muted}#3}}}\par\smallskip}
\newtcolorbox{paperbox}[2]{enhanced,breakable,colback=paperbg,colframe=paperbg,
  boxrule=0pt,leftrule=2.4pt,arc=0pt,outer arc=0pt,
  borderline west={2.4pt}{0pt}{paperrule},
  left=9pt,right=9pt,top=6pt,bottom=7pt,before skip=8pt,after skip=10pt,
  halign=justify,before upper={\boxlabel{paperlab}{#1}{#2}\PaperMacros\addfontfeatures{Ligatures=TeX}\setlength{\parskip}{3pt}\hyphenpenalty=50 \exhyphenpenalty=50 \justifying}}
\newtcolorbox{texverbbox}[2]{enhanced,breakable,colback=paperbg,colframe=paperbg,
  boxrule=0pt,borderline west={2.4pt}{0pt}{paperrule},arc=0pt,outer arc=0pt,
  left=9pt,right=9pt,top=6pt,bottom=6pt,before skip=8pt,after skip=10pt,
  before upper={\boxlabel{paperlab}{#1}{#2}}}
\newtcolorbox{leanstmt}[2]{enhanced,breakable,colback=leanbg,colframe=leanbg,
  boxrule=0pt,borderline west={2.4pt}{0pt}{leanrule},arc=0pt,outer arc=0pt,
  left=9pt,right=6pt,top=6pt,bottom=6pt,before skip=8pt,after skip=12pt,
  before upper={\boxlabel{leanrule}{#1}{#2}}}
\newtcolorbox{leandef}[1]{enhanced,breakable,colback=defbg,colframe=defframe,
  boxrule=0.4pt,arc=1.5pt,outer arc=1.5pt,
  left=9pt,right=6pt,top=4pt,bottom=4pt,before skip=6pt,after skip=8pt}
\newtcolorbox{leanmisc}[2]{enhanced,breakable,colback=miscbg,colframe=miscbg,
  boxrule=0pt,borderline west={2.4pt}{0pt}{faint},arc=0pt,outer arc=0pt,
  left=9pt,right=6pt,top=5pt,bottom=5pt,before skip=8pt,after skip=10pt,
  before upper={\boxlabel{muted}{#1}{#2}}}
\newtcolorbox{notebox}{enhanced,breakable,colback=white,colframe=defframe,boxrule=0.5pt,arc=2pt,
  left=10pt,right=10pt,top=8pt,bottom=8pt,fontupper=\sffamily\small,halign=flush left}

% ---------- coverage chips ----------
\DeclareRobustCommand\chip[3]{{\setlength{\fboxsep}{1.6pt}\colorbox{#1}{\sffamily\fontsize{7}{8}\selectfont\textcolor{#2}{\strut\,#3\,}}}}
\DeclareRobustCommand\covchip[2]{\ifstrequal{#1}{green}{\chip{chipgreenbg}{chipgreen}{#2}}{%
  \ifstrequal{#1}{amber}{\chip{chipamberbg}{chipamber}{#2}}{\chip{chipbluebg}{chipblue}{#2}}}}

% ---------- equation keys / refs in paper excerpts ----------
\newcommand\eqkey[1]{{\normalfont\sffamily\fontsize{7.5}{9}\selectfont\color{paperlab}(#1)}}
\newcommand\refkey[1]{{\normalfont\sffamily\fontsize{8}{9}\selectfont\color{paperlab}#1}}

% ---------- theorem environments for the excerpts (amsart-like) ----------
\theoremstyle{plain}
\newtheorem{theorem}{Theorem}
\newtheorem{lemma}[theorem]{Lemma}
\newtheorem{proposition}[theorem]{Proposition}
\newtheorem{corollary}[theorem]{Corollary}
\newtheorem{definition}[theorem]{Definition}
\newtheorem{conjecture}[theorem]{Conjecture}
\renewcommand{\labelenumi}{(\theenumi)}

% ---------- headings ----------
\newcommand\partrule{{\color{defframe}\rule{\linewidth}{0.6pt}}}
\newcommand\stmthead[6]{% #1 title #2 key #3 anchor #4 chip #5 srcline #6 pdf string
  \clearpage
  \phantomsection\hypertarget{#3}{}%
  \addcontentsline{toc}{section}{\texorpdfstring{\tocstmt{#1}{#2}{#4}}{#6}}%
  \markboth{#6}{#6}%
  {\noindent\sffamily\color{accent}{\fontsize{19}{23}\selectfont\bfseries #1}%
   \ifstrempty{#2}{}{\hspace{0.7em}{\ttfamily\large\color{muted}#2}}%
   \hfill{\fontsize{7.5}{9}\selectfont\color{faint}Statements.lean:#5}\par}%
  \vspace{2pt}{\noindent\ifstrempty{#4}{}{#4}}\par
  \vspace{-2pt}{\noindent\partrule}\par\vspace{4pt}}
\newcommand\parthead[5]{% #1 title #2 anchor #3 srcline #4 pdf string #5 newpage(1/0)
  \ifstrequal{#5}{1}{\clearpage}{\Needspace{8\baselineskip}\vspace{10pt}}%
  \phantomsection\hypertarget{#2}{}%
  \addcontentsline{toc}{section}{\texorpdfstring{\tocpart{#1}}{#4}}%
  \markboth{#4}{#4}%
  {\noindent\sffamily\color{accent}{\fontsize{15}{19}\selectfont\bfseries #1}%
   \hfill{\fontsize{7.5}{9}\selectfont\color{faint}Statements.lean:#3}\par}%
  \vspace{-1pt}{\noindent\partrule}\par\vspace{3pt}}
\newcommand\defhead[5]{% #1 id #2 name(latex) #3 anchor #4 srcline #5 pdf string
  \Needspace{7\baselineskip}\vspace{9pt}%
  \phantomsection\hypertarget{#3}{}\bookmark[dest=#3,level=2]{#5}%
  {\noindent\colorbox{accent}{\sffamily\fontsize{7.5}{9}\selectfont\bfseries\color{white}\strut #1}%
   \hspace{0.55em}{\ttfamily\bfseries\color{accent}#2}%
   \hfill{\sffamily\fontsize{7.5}{9}\selectfont\color{faint}Statements.lean:#4}\par}\vspace{1pt}}
\newcommand\glosshead[5]{% #1 id #2 rest #3 anchor #4 srcline #5 pdf
  \Needspace{6\baselineskip}\vspace{8pt}%
  \phantomsection\hypertarget{#3}{}\bookmark[dest=#3,level=2]{#5}%
  {\noindent{\sffamily\bfseries\color{accent}#1}\hspace{0.5em}{\sffamily\color{muted}—}\hspace{0.5em}{\color{accent}\bfseries #2}%
   \hfill{\sffamily\fontsize{7.5}{9}\selectfont\color{faint}Statements.lean:#4}\par}\vspace{1pt}}
\newcommand\ctxhead[5]{% #1 title #2 key #3 anchor #4 srcline #5 pdf
  \Needspace{6\baselineskip}\vspace{10pt}%
  \phantomsection\hypertarget{#3}{}\bookmark[dest=#3,level=2]{#5}%
  {\noindent{\sffamily\fontsize{7.2}{9}\selectfont\bfseries\color{muted}PAPER CONTEXT}%
   \hspace{0.6em}{\ttfamily\color{muted}#2}\hfill{\sffamily\fontsize{7.5}{9}\selectfont\color{faint}Statements.lean:#4}\par
   \noindent{\sffamily\large\bfseries\color{accent}#1}\par}\vspace{2pt}}
\newcommand\subhead[4]{% #1 title #2 anchor #3 srcline #4 pdf
  \Needspace{5\baselineskip}\vspace{8pt}%
  \phantomsection\hypertarget{#2}{}\bookmark[dest=#2,level=2]{#4}%
  {\noindent{\sffamily\bfseries\color{accent}#1}\hfill{\sffamily\fontsize{7.5}{9}\selectfont\color{faint}Statements.lean:#3}\par}}
\DeclareRobustCommand\tocstmt[3]{#1\hspace{0.6em}{\ttfamily\small\color{muted}#2}\hspace{0.6em}#3}
\DeclareRobustCommand\tocpart[1]{{\bfseries #1}}

% ---------- table of contents ----------
\makeatletter
\renewcommand\l@section[2]{%
  \par\noindent{\sffamily\small #1\nobreak\leaders\hbox{$\m@th\mkern 4mu\hbox{\color{faint}.}\mkern 4mu$}\hfill\nobreak\hbox to 2em{\hss #2}}\par\vspace{1.5pt}}
\makeatother

% ---------- header / footer ----------
\pagestyle{fancy}
\fancyhf{}
\renewcommand{\headrulewidth}{0.3pt}
\renewcommand{\headrule}{\ifdim\headrulewidth>0pt{\color{defframe}\hrule height \headrulewidth}\fi}
\fancyhead[L]{\sffamily\fontsize{7.5}{9}\selectfont\color{muted}Statements.lean \textperiodcentered{} paper statements and their Lean encodings}
\fancyhead[R]{\sffamily\fontsize{7.5}{9}\selectfont\color{accent}\rightmark}
\fancyfoot[C]{\sffamily\fontsize{7.5}{9}\selectfont\color{muted}\thepage\ /\ \pageref*{LastPage}}
\fancypagestyle{plain}{\fancyhf{}\renewcommand{\headrulewidth}{0pt}\fancyfoot[C]{\sffamily\fontsize{7.5}{9}\selectfont\color{muted}\thepage\ /\ \pageref*{LastPage}}}
